% GENERATED FILE. Edit canonical lessons, then run npm run generate:books.
% canonical-source-hash: fnv1a64:2e1073dbad0e20ad
% canonical-lessons: MW-C09-hear-patni, MW-C09-patni, MW-C09-hear-bachcha, MW-W09-cha, MW-C09-bachcha, MW-R09-family-new-five, MW-C09-family-twelve, MW-R09-script-close

\chapter{Wife and Child: One New Sign}
\label{ch:mw-family-twelve}
\hlchaptermodality{13}

\begin{chapteropening}
\textbf{By the end of this chapter:} \emph{I can add wife and child to a twelve-label family map while recognizing that the source records an alternative wife label.}

Wife and child begin by ear, add only \mw{च}, and join the earlier labels in a non-compensatory four-skill payoff.
\end{chapteropening}

\section[patnī]{Hear \emph{patnī} — wife}

\label{lesson:MW-C09-hear-patni}



\begin{quote}
Say husband and maternal grandmother, write \textbf{\mw{सा}}, then recall the last three-word four-skill checkpoint.
\end{quote}

\subsection*{What to know first}
\begin{quote}
\emph{patnī} — \textbf{wife}
\end{quote}

The same source also gives \emph{bāī sā}. That documented alternative is not an error; this step secures one form at a time. Writing still waits.

\begin{tcolorbox}[breakable,colback=teal!4,colframe=teal!35!black,title={Before you move on}]
Source: \href{https://www.marwaripathshala.com/marwari-lesson-4-english}{Marwari Pathshala, Lesson 4}.
\end{tcolorbox}
\section[patnī]{\mw{पत्नी} — known signs, one careful join}

\label{lesson:MW-C09-patni}



\begin{quote}
Say \emph{patnī}. Write \textbf{\mw{पति}}, \textbf{\mw{्}}, \textbf{\mw{न}}, and \textbf{\mw{ी}} separately; write \textbf{\mw{नानी}}, then say the respectful \textbf{\mw{राम} \mw{राम} \mw{सा}} greeting.
\end{quote}

\subsection*{What to know first}
\begin{quote}
\textbf{\mw{प}} + \textbf{\mw{त्}} + \textbf{\mw{नी}} $\to$ \textbf{\mw{पत्नी}} — \emph{patnī} — wife
\end{quote}

\subsection*{Your turn}
Copy once. Cover it, wait five seconds, then write it. Check the virama join before checking the final long-ī sign.

\begin{tcolorbox}[breakable,colback=teal!4,colframe=teal!35!black,title={Before you move on}]
Source: \href{https://www.marwaripathshala.com/marwari-lesson-4-english}{Marwari Pathshala, Lesson 4}.
\end{tcolorbox}
\section[bachchā]{Hear \emph{bachchā} — child}

\label{lesson:MW-C09-hear-bachcha}



\begin{quote}
Say wife, husband, and maternal grandfather, complete the familiar greeting exchange, then write \textbf{\mw{दादा}}.
\end{quote}

\subsection*{What to know first}
\begin{quote}
\emph{bachchā} — \textbf{child}
\end{quote}

Hear the held middle consonant. The label does not specify a child's gender; writing waits until the new consonant has a separate, small step.

\begin{tcolorbox}[breakable,colback=teal!4,colframe=teal!35!black,title={Before you move on}]
Source: \href{https://www.marwaripathshala.com/marwari-lesson-4-english}{Marwari Pathshala, Lesson 4}.
\end{tcolorbox}
\section[cha]{\mw{च} — unaspirated beside familiar \mw{छ}}

\label{lesson:MW-W09-cha}



\begin{quote}
Say child and wife. Write familiar \textbf{\mw{छ}} and independent \textbf{\mw{आ}}; today's sound has no following breath.
\end{quote}

\subsection*{Script — one sign}
\begin{quote}
\textbf{\mw{च}} — \emph{cha}
\end{quote}

\subsection*{Your turn}
Trace once, copy once, cover the model, and write it once from sound.

\begin{tcolorbox}[breakable,colback=teal!4,colframe=teal!35!black,title={Before you move on}]
Source: \href{https://www.unicode.org/charts/PDF/U0900.pdf}{Unicode Devanagari chart}.
\end{tcolorbox}
\section[bachchā]{\mw{बच्चा} — child on the page}

\label{lesson:MW-C09-bachcha}



\begin{quote}
Say child and wife. Write \textbf{\mw{ब}}, \textbf{\mw{च}}, \textbf{\mw{्}}, \textbf{\mw{ा}}, familiar \textbf{\mw{भ}}, then write \textbf{\mw{पत्नी}} and \textbf{\mw{दादी}} from memory.
\end{quote}

\subsection*{What to know first}
\begin{quote}
\textbf{\mw{ब}} + \textbf{\mw{च्}} + \textbf{\mw{चा}} $\to$ \textbf{\mw{बच्चा}} — \emph{bachchā} — child
\end{quote}

\subsection*{Your turn}
Copy once. Cover it, wait five seconds, and write it. Check the doubled middle before checking the final vowel mark.

\begin{tcolorbox}[breakable,colback=teal!4,colframe=teal!35!black,title={Before you move on}]
Source: \href{https://www.marwaripathshala.com/marwari-lesson-4-english}{Marwari Pathshala, Lesson 4}.
\end{tcolorbox}
\section[Practice]{Retrieve the five new relationship labels}

\label{lesson:MW-R09-family-new-five}



\begin{quote}
Write \textbf{\mw{त}}, \textbf{\mw{च}}, and \textbf{\mw{आभार}}, give the gratitude meaning, then recall the earlier three-word payoff.
\end{quote}

\subsection*{Your turn}
Hear all five in a mixed order, say each from meaning, read the five cards, then write the two cards named aloud. Repair only a missed card.

\begin{tcolorbox}[breakable,colback=teal!4,colframe=teal!35!black,title={Before you move on}]
Source: \href{https://www.marwaripathshala.com/marwari-lesson-4-english}{Marwari Pathshala, Lesson 4}.
\end{tcolorbox}
\section[Practice]{Payoff — a twelve-label family map}

\label{lesson:MW-C09-family-twelve}



\begin{quote}
Recall the seven-label and three-label checkpoints, then give the formal gratitude response. These are vocabulary options, not a prescribed family structure.
\end{quote}

\subsection*{Your turn}
1. \textbf{Listen:} identify six randomly heard labels. 2. \textbf{Speak:} produce six different labels from relationship cues. 3. \textbf{Read:} match all twelve printed labels to meanings. 4. \textbf{Write:} write four heard labels without a model.

Score each skill separately. A strong spoken score cannot compensate for a missed reading or writing score.

\begin{tcolorbox}[breakable,colback=teal!4,colframe=teal!35!black,title={Before you move on}]
Source: \href{https://www.marwaripathshala.com/marwari-lesson-4-english}{Marwari Pathshala, Lesson 4}.
\end{tcolorbox}
\section[Practice]{Close the chapter — two signs, two words}

\label{lesson:MW-R09-script-close}



\begin{quote}
Write earlier \textbf{\mw{द}}, \textbf{\mw{व}}, \textbf{\mw{ह}}, and \textbf{\mw{परिवार}}, then name the four skills in the twelve-label payoff.
\end{quote}

\subsection*{Your turn}
From sound alone, write \textbf{\mw{त}}, \textbf{\mw{च}}, \textbf{\mw{पत्नी}}, and \textbf{\mw{बच्चा}}. Check each new sign before checking the complete word.

\begin{tcolorbox}[breakable,colback=teal!4,colframe=teal!35!black,title={Before you move on}]
Source: \href{https://www.unicode.org/charts/PDF/U0900.pdf}{Unicode Devanagari chart}.
\end{tcolorbox}
